\documentclass[a4paper,11pt]{article}

\usepackage[T1]{fontenc}
\usepackage[utf8]{inputenc}
\usepackage[italian]{babel}
\usepackage[margin=2.5cm]{geometry}
\usepackage{microtype}
\usepackage{setspace}
\usepackage{graphicx}
\usepackage{float}
\usepackage{booktabs}
\usepackage{multirow}
\usepackage{array}
\usepackage{tabularx}
\usepackage{longtable}
\usepackage{enumitem}
\usepackage{caption}
\usepackage{subcaption}
\usepackage{xcolor}
\usepackage{hyperref}
\usepackage{amsmath,amssymb,amsthm,mathtools}
\usepackage{bm}
\usepackage{siunitx}
\usepackage{upgreek}
\usepackage{dsfont}
\usepackage{nicefrac}
\usepackage{algorithm}
\usepackage{algpseudocode}
\usepackage{listings}
\usepackage[numbers,sort&compress]{natbib}

\onehalfspacing
\setlength{\parindent}{0pt}
\setlength{\parskip}{0.5em}

\hypersetup{
    colorlinks=true,
    linkcolor=blue,
    citecolor=blue,
    urlcolor=blue,
    pdfauthor={Autore},
    pdftitle={Titolo del lavoro}
}

\newtheorem{theorem}{Teorema}[section]
\newtheorem{lemma}[theorem]{Lemma}
\newtheorem{proposition}[theorem]{Proposizione}
\newtheorem{definition}[theorem]{Definizione}
\newtheorem{remark}[theorem]{Osservazione}
\newtheorem{corollary}[theorem]{Corollario}

\DeclareMathOperator*{\argmin}{arg\,min}
\DeclareMathOperator*{\argmax}{arg\,max}
\newcommand{\trace}{\operatorname{Tr}}

\newcommand{\R}{\mathbb{R}}
\newcommand{\N}{\mathbb{N}}
\newcommand{\C}{\mathbb{C}}
\newcommand{\E}{\mathbb{E}}
\newcommand{\Var}{\operatorname{Var}}
\newcommand{\Cov}{\operatorname{Cov}}
\newcommand{\vect}[1]{\mathbf{#1}}
\newcommand{\mat}[1]{\mathbf{#1}}
\newcommand{\norm}[1]{\left\lVert #1 \right\rVert}
\newcommand{\abs}[1]{\left\lvert #1 \right\rvert}

\title{Titolo del lavoro}
\author{Nome Cognome \\[0.3em]
\small{Dipartimento / Universit\`a / Team} }
\date{\today}

\begin{document}

\maketitle

\begin{abstract}
Questo documento rappresenta un template di articolo scientifico in LaTeX con impostazione tipografica pulita e leggibile.\`E progettato per supportare la suddivisione del contenuto in file separati, cos\`i da poter organizzare introduzione, metodi, risultati e appendici in sezioni indipendenti.\
\end{abstract}

\tableofcontents
\newpage

\section{Introduzione}
Il progetto Prosody To Music nasce dall'idea di associare una composizione musicale a testo in lingua italiana.
Spesso, durante la lettura, può succedere di seguire un flusso melodico spontaneo, dando intonazione alle parole; difatti, la lingua italiana è di 
per sé fortemente musicale e ritmica.
Sfruttando dunque la prosodia ricavabile dal susseguirsi di parole, questo programma è in grado di produrre una base musicale 
coerente con l'input testuale.


\section{Obiettivi}
Il progetto ha come obiettivo principale la definizione di una pipeline coerente per trasformare un input testuale in una rappresentazione musicale espressiva 
e controllabile.
Nello specifico, il progetto mira al raggiungimento dei seguenti obiettivi:

\begin{enumerate}
    \item Prendere in input un testo ed estrapolare, secondo le regole della lingua italiana, varie tipologie di accenti, pause respiratorie e 
    significati sintattici.
    \item Produrre uno schema ritmico musicale coerente con la distribuzione degli accenti ritmici e dei silenzi.
    \item Generare una melodia coerente con il ritmo ottenuto precedentemente, lasciando libertà di personalizzazione all'utente.
\end{enumerate}
\section{Metodologia in breve}
Al fine di raggiungere gli obiettivi presentati sarà seguita la seguente pipeline:
\begin{enumerate}
    \item \textbf{Analisi del testo:} l'analisi del testo sarà svolta parallelamente dal punto di vista semantico e da quello ritmico per completare un processo di analisi prosodica del testo.  La semantica servirà a generare un'atmosfera musicale coerente con il significato e il contesto emotivo del testo. La prosodia servirà ad individuare gli accenti ritmici e le pause per la produzione dello schema ritmico.
    \item \textbf{Produzione dello schema ritmico:} traduzione dello studio prosodico in una struttura ritmicas coerente con il testo di input.
    \item \textbf{Addestramento del modello:} per la generazione della melodia, in questa fase si addestrerà un modello Transformer su un dataset di brani musicali.
    \item \textbf{Generazione musicale:} sarà generato un file MIDI coerente con il ritmo e con la semantica del testo, attraverso il Transformer precedentemente allenato.
    \item \textbf{Sintesi del brano:} il file MIDI sarà infine sintetizzato e convertito in un file audio (WAVE).      
\end{enumerate}


\section{Specifiche del progetto}
Il linguaggio di programmazione adottato per la realizzazione di questo progetto è Python, in quanto dotato di numerose librerie estremamte utili 
per l'implementazione. Inoltre, per l'addestramento dei modelli sono previsti dataset reperibili online contenenti dati e strutturati e pertinenti.
I dataset e le librerie saranno presentate nel dettaglio del capitolo successivo.

\section{Dal testo agli accenti al ritmo}
\section{Approccio generale}
Come prima cosa si è deciso di affrontare il problema in maniera del rilevamento degli accenti a partire da un testo in lingua italiana.
Per riuscire ad arrivare a un modello capace di riuscire nel compito dell'induviazione degli accenti è necessario importaqre dei dataset importati per il task. Fatto ciò tali dataset saranno usati per creare per appunto il modello che data una generica frase in italiano possa individuare tutti gli accenti con correttezza
\subsection{Q2Stress}
Q2Stress è un dataset di frasi in lingua italiana, ognuna delle quali è annotata con informazioni sugli accenti presenti. Questo dataset è stato utilizzato per addestrare il modello di riconoscimento degli accenti, fornendo esempi concreti di come gli accenti si manifestano nel linguaggio naturale.
\subsection{phonItaliaR}
phonItaliaR è un dataset di frasi in lingua italiana, ognuna delle quali è annotata con informazioni sulle pronunce e sugli accenti presenti. Questo dataset è stato utilizzato per addestrare il modello di riconoscimento degli accenti, fornendo esempi concreti di come gli accenti si manifestano nel linguaggio naturale.
\section{Definizione del modello}
\begin{definition}
Un modello matematico \`e una rappresentazione semplificata di un fenomeno reale, costruita per descrivere relazioni quantificabili tra variabili.
\end{definition}

Per esempio, una funzione di costo pu\`o essere definita come
\begin{equation}
\mathcal{L}(\theta) = \frac{1}{n} \sum_{i=1}^{n} \ell(y_i, f_\theta(x_i)) + \lambda \|\theta\|_2^2.
\end{equation}

\section{Ottimizzazione}
L'ottimizzazione di questa espressione richiede tecniche numeriche, come gradient descent o metodi quasi-Newton, a seconda della complessit\`a del problema e della dimensione dello spazio dei parametri.

\begin{remark}
La scelta della funzione di perdita influenza in modo diretto la robustezza del modello rispetto a rumore, outlier e distribuzioni non gaussiane.
\end{remark}


\section{Selezione degli strumenti}
\section{Metodologie utilizzate}
Una volta completata l'analisi prosodica ed emotiva si può passare alla parte di produzione mulsicale.

Per svolgere questo task si farà uso di un transformer

\begin{table}[H]
\centering
\caption{Esempio di tabella sintetica.}
\begin{tabular}{lcc}
\toprule
Metodo & Valore medio & Deviazione standard \\
\midrule
A & 2.34 & 0.41 \\
B & 2.85 & 0.27 \\
C & 3.12 & 0.33 \\
\bottomrule
\end{tabular}
\end{table}

\section{Stima del modello}
Un risultato di interesse \`e dato dalla seguente stima:
\begin{equation}
\hat{\beta} = \argmin_{\beta \in \R^p} \left\{ \|\vect{y} - \mat{X}\beta\|_2^2 + \lambda \|\beta\|_1 \right\}.
\end{equation}

Questo tipo di formulazione evidenzia come la regolarizzazione possa stabilire un compromesso tra accuratezza del fit e generalizzazione del modello.


\section{Conclusioni}
\section{Valutazione finale}
Il risultato ottentuto risulta soddisfaciente, il programma sviluppato riesce a raggiungere gli obiettivi preposti.

Il programma riesce a generare della musica coerente con il un prompt testuale di input, inoltre il ritmo della canzone rispecchia fedelmente la distrubuzione degli accenti delle parole.



\section{Sviluppi futuri}
Questo progetto lascia la porta aperta a numerosi sviluppi futuri:
\begin{itemize}
    \item Ampliare i dataset di addestramento dei modelli di Machine Learning e Deep Learning utilizzati.
    \item Offrire l'estensione ad altre lingue al di fuori dell'italilano, come l'inglese o lo spagnolo, per rendere il sistema più versatile.
    \item Possibilità di aggiungere la generazione di una voce che canta le parole del testo riproducendo fedelmente la prosodia.
    \item Aggiungere una memoria persistente per evitare di ripetere ad ogni esecuzione l'addestramento del Random Forest nella fase di analisi prosodica. 
\end{itemize}




\appendix
\section{Appendice A}
\section*{A.1 Esempio di appendice}
In appendice si possono inserire dettagli secondari, dimostrazioni o dati supplementari. Ad esempio, la seguente identit\`a \`e utile in numerosi contesti:
\begin{equation}
\sum_{k=0}^{n} \binom{n}{k} = 2^n.
\end{equation}

Questa forma pu\`o essere usata per riportare risultati accessori senza interrompere il flusso principale dell'articolo.


% Se si desidera usare una bibliografia esterna, decommentare le righe seguenti
% \bibliographystyle{plainnat}
% \bibliography{bibliografia}

\end{document}
